% ભાગ: ગણો-વિધેયો-સંબંધો
% પ્રકરણ: અંકગણિતીકરણ
% વિભાગ: પ્રતિચિંતનો

\documentclass[../../../include/open-logic-section]{subfiles}

\begin{document}

\olfileid{sfr}{arith}{ref}
\olsection{કેટલાક તાત્ત્વિક પ્રતિચિંતનો}

પ્રાવિધિક વિગતો આટલી થઈ. પરંતુ તેનાથી સિદ્ધ શું થયું?

એવું લગભગ નિર્વિવાદ છે કે તેનાથી આપણને થોડું {સુંદર} શુદ્ધ ગણિત
મળ્યું. ઉપરાંત કેટલીક ગહન વૈચારિક સિદ્ધિઓ મળી. પૂર્ણતા ગુણધર્મ
વાસ્તવિક અને સંમેય સંખ્યાઓ વચ્ચેનો નિર્ણાયક તફાવત વ્યક્ત કરે છે એવું
જોવું એક ગહન અંતર્દૃષ્ટિ હતું. વધુમાં, ડેડેકિન્ડ કાપ તરીકે વાસ્તવિક
સંખ્યાઓની સ્પષ્ટ રચના વિશ્લેષણના વિષયને મજબૂત પાયો આપે છે. કાપો
\emph{પૂર્ણ ક્રમિત ક્ષેત્ર} બનાવે છે, તેથી આ ખ્યાલ સુસંગત છે તે આપણે જાણીએ છીએ.

આ બધું સ્વીકાર્યા પછી પણ આ સિદ્ધિઓ વિશે કેટલીક શંકાઓ જણાવવી જોઈએ.

પહેલી વાત, વાસ્તવિક સંખ્યાઓને કાપો તરીકે વિચારવી એ તેમના પરિચિત
(કદાચ અનંત) દશાંશ વિસ્તરણોના રૂપમાં વિચારવા કરતાં \emph{વધુ} ચુસ્ત
છે કે નહીં તે સ્પષ્ટ નથી. વાસ્તવિક સંખ્યાઓની આ બીજી ``રચના'' કોશી
શ્રેણી વડે થતી રચના સાથે થોડું સામ્ય ધરાવે છે; પરંતુ હકીકતમાં
સત્તરમી સદીના આરંભથી ગણિતજ્ઞો તેને મૂળભૂત રીતે જાણતા હતા
(\olref[cauchy]{sec} જુઓ). ચુસ્તતામાં ખરેખરનો વધારો વાસ્તવિક સંખ્યાઓ
પૂર્ણતા ગુણધર્મ ધરાવે છે એ સમજવાથી થયો. ચોક્કસ ગણો તરીકે વાસ્તવિક
સંખ્યાઓ રચવાની ક્ષમતા કદાચ પોતાની રીતે એટલી રસપ્રદ નથી.

પૂર્ણાંકોનું કે સંમેય સંખ્યાઓનું (ઘણું સરળ) અંકગણિતીકરણ તે ક્ષેત્રોમાં
ચુસ્તતા વધારે છે કે નહીં તે તો હજી ઓછું સ્પષ્ટ છે. અહીં એક સરળ અવલોકન
ઉપયોગી છે. પ્રાકૃતિક સંખ્યાઓની ક્રમયુક્ત જોડોના સામ્ય વર્ગો તરીકે
પૂર્ણાંકોની \emph{રચના} કર્યા પછી, પૂર્ણાંકોની ક્રમયુક્ત જોડોના સામ્ય
વર્ગો તરીકે સંમેય સંખ્યાઓની રચના કર્યા પછી અને સંમેય સંખ્યાઓના ગણો
તરીકે વાસ્તવિક સંખ્યાઓ રચ્યા પછી આપણે તરત જ એ રચનાઓને \emph{ભૂલી
જઈએ છીએ}. ખાસ કરીને, ગણિતીય સાબિતીમાં કોઈ આ રચનાઓને ક્યારેય
\emph{ઉપયોગમાં લેવા} ઇચ્છશે નહીં (અલબત્ત, રચનાઓનો વ્યવહાર ધાર્યા મુજબ
છે એવી સાબિતી સિવાય). પ્રાકૃતિક સંખ્યાઓના ગણોના ગણોના ગણોના ગણોના
ગણોના ગણોના કોઈ ગણ વિશે બોલવા કરતાં સીધું કોઈ વાસ્તવિક સંખ્યા વિશે
બોલવું ઘણું સરળ છે.
%(કારણ કે વાસ્તવિક સંખ્યાઓ સંમેય સંખ્યાઓના ગણો તરીકે રચાઈ; અને સંમેય સંખ્યાઓ પૂર્ણાંકોની ક્રમયુક્ત જોડોના સામ્ય વર્ગો, એટલે કે પૂર્ણાંકોના ગણોના ગણોના ગણો તરીકે રચાઈ; વગેરે.)
%2/3 = [2,3]
%  એટલે કે {<2,3>, ... }
%  { {{2},{2,3}}, ... }
%
%વાસ્તવિક સંખ્યાઓ છે
%  સંમેય સંખ્યાઓના ગણો
%
%સંમેય સંખ્યાઓ છે
%  પૂર્ણાંકોના ગણોના ગણોના ગણો
%
%પૂર્ણાંકો છે
%  પ્રાકૃતિક સંખ્યાઓના ગણોના ગણોના ગણો

આ વ્યાખ્યાઓ આપણને પૂર્ણાંકો, સંમેય સંખ્યાઓ કે વાસ્તવિક સંખ્યાઓ
\emph{અધિભૌતિક અર્થમાં શું છે} તે કહે છે એવો દાવો સૌથી વધુ
શંકાસ્પદ છે. એટલે કે, વાસ્તવિક સંખ્યાઓ (ધારો કે) ચોક્કસ ગણો
(ગણોના ગણોના ગણો\ldots) જ \emph{છે} એ શંકાસ્પદ છે. આવા મતનો મુખ્ય
અવરોધ એ છે કે રચના ઘણી જુદી રીતે થઈ શકી હોત. વાસ્તવિક સંખ્યાઓના
કિસ્સામાં ખરેખર રસપ્રદ રીતે જુદી રચનાઓ છે (\olref[cauchy]{sec} જુઓ).
પરંતુ તદ્દન જુદી રચનાઓ મેળવવાની આ એક નજીવી રીત છે: જેમ
\olref[sfr][rel][ref]{sec}માં જોયું, ક્રમયુક્ત જોડની વ્યાખ્યા સહેજ જુદી
રીતે કરી શકાઈ હોત; આ વૈકલ્પિક ખ્યાલ વાપર્યો હોત તો આપણી રચનાઓ એટલી
જ સારી રીતે કામ કરત, પરંતુ અંતે મળતી વસ્તુઓ જુદી હોત. આમ પૂર્ણાંકો,
સંમેય સંખ્યાઓ અને વાસ્તવિક સંખ્યાઓની ઘણી હરીફ ગણસૈદ્ધાંતિક રચનાઓ છે.
હવે પૂર્ણાંકો (ધારો કે) \emph{આ} ગણો છે, \emph{તે} નહીં, એવો દાવો
મનસ્વી (અને મૂંઝવણજનક) થશે. (\olref[sfr][rel][ref]{sec}ની જેમ આ પણ
\citealt{Benacerraf1965}એ પ્રસિદ્ધ કરેલી દલીલનું એક દૃષ્ટાંત છે.)

બીજો મુદ્દો પણ ઉઠાવવો જોઈએ: આપણી રચનાઓમાં કશુંક બહુ \emph{વિચિત્ર}
છે. આપણે પ્રાકૃતિક સંખ્યાઓથી શરૂ કર્યું. પછી પૂર્ણાંકો રચ્યા અને
``પૂર્ણાંકોનું $0$'', એટલે કે $\equivrep{0,0}{\Intequiv}$ રચ્યું.
પરંતુ $0 \neq \equivrep{0,0}{\Intequiv}$. ખરેખર, આપણી રચનાઓ મુજબ
\emph{કોઈ પણ} પ્રાકૃતિક સંખ્યા પૂર્ણાંક નથી. આ ખૂબ જ પ્રતિઅંતઃપ્રજ્ઞ
લાગે છે. વળી \olref[sfr][set][imp]{sec}માં આપણે બહુ દલીલ કર્યા વિના
$\Nat \subseteq \Rat$ એવો દાવો કર્યો હતો. જો આ રચનાઓ સંખ્યાઓ
\emph{બરાબર શું છે} તે કહેતી હોય, તો એ દાવો સહેજે ખોટો હતો.

હવે થોડું દૂરથી જોઈએ તો આપણે ક્યાં પહોંચ્યા? સહજ ગણસિદ્ધાંતમાં કામ
કરીને અને પ્રાકૃતિક સંખ્યાઓને આપેલી માનીને, પૂર્ણાંકો, સંમેય સંખ્યાઓ
અને વાસ્તવિક સંખ્યાઓને ચોક્કસ ગણો \emph{તરીકે} ગણી શકીએ છીએ. આ અર્થમાં
આ વસ્તુઓના સિદ્ધાંતોને ગણસિદ્ધાંતમાં \emph{સ્થાપિત} કરી શકીએ છીએ.
પરંતુ આ સ્થાપનનું તાત્ત્વિક મહત્ત્વ એટલું સીધું નથી.

અલબત્ત, આમાંથી કંઈ અંતિમ શબ્દ નથી!{} મુદ્દો માત્ર આ છે: વાસ્તવિક
સંખ્યાઓનું અંકગણિતીકરણ ગહન તાત્ત્વિક મહત્ત્વ \emph{ધરાવે છે} એવું
બતાવવા માટે કેટલીક વધારાની \emph{તાત્ત્વિક} દલીલ જરૂરી થશે.

%વધુમાં: આ આખા પ્રકરણ દરમિયાન આપણે પ્રાકૃતિક સંખ્યાઓને સીધી
%\emph{આપેલી} માની છે. પરંતુ તેમના વિશે આપણે શું કરી શકીએ?
%આગળનું પ્રકરણ એ જ વિષય લે છે.

\end{document}
